% ---
% id: meat-core/rem:tightness-by-construction
% kind: rem
% title: "Tightness read off the construction of the frame"
% summary: >-
%   Deciding tightness of a given τ_A is a rank decision on a computed
%   spectrum; the diagonal frame exhibits it by construction, the full frame
%   exhibits nothing. Read at a positive window, the diagonality test would
%   admit the frame-free rows on a frame that is not tight.
% prominence: standard
% status: accepted
% uses: [def:support-frame, lem:choi-marginal, thm:choi, prop:frame-free-rows]
% mentions: []
% assumes: []
% refs: []
% provenance: newly-derived
% lean:
% ---

\begin{remark}\label{rem:tightness-by-construction}
Tightness of a support frame is $\ker W_{\tau}=0$, equivalently
$d_{A_0}=\operatorname{rank}\tau_A$ (Lemma~\ref{lem:choi-marginal}\,(iii)),
so establishing it \emph{of a given} $\tau_A$ is a rank decision on a
computed spectrum: the decision that \eqref{eq:iota} removes, since
Lemma~\ref{lem:choi-marginal} and Theorem~\ref{thm:choi} hold on every
support frame. What can be known without it is whether the construction
that produced the frame exhibits $\ker W_{\tau}=0$ without computing a rank.

Of the two constructions of Definition~\ref{def:support-frame}, the
diagonal frame~(b) does: there
$\operatorname{ran}\iota=\operatorname{supp}\tau_A$ holds by construction,
the index set $\{i:t_i\neq0\}$ being fixed by a bit-level comparison of
declared binary64 data, which is exact because~(b) is available only for a
$\tau_A$ diagonal in the declared basis, where its diagonal entries
\emph{are} its eigenvalues. This covers the case where every $t_i\neq0$
and construction~(b) returns $\iota=\idon{A}$. The full frame taken as
construction~(a) is built from $\idon{A}$ and from no property of $\tau_A$
at all, and therefore exhibits nothing about $\operatorname{rank}\tau_A$,
however nonsingular $\tau_A$ looks.

The diagonality on which~(b) rests is itself decided by comparing declared
entries with zero. Read at a positive window instead, a singular $\tau_A$
with small off-diagonal entries and every $t_i\ne0$ would pass as
diagonal, and~(b) would return the full frame, which is not tight. The
frame-free form of Proposition~\ref{prop:frame-free-rows} read off that
construction would then be imposed on a non-tight frame, where it imposes
the row on the directions of $\ker W_{\tau}$ as well, cutting $\Feas^{J}$
down to a subset, in general proper, whose image under $\chi_{\Qsys}$ need
not be all of $\Feas$; by Proposition~\ref{prop:frame-free-rows}\,(ii) an
infimum over that subset can only come out higher, the unsafe direction for
a key length. For a party-local row of the general format, the
unconjugated form $Y_k\bowtie_k0$ of
Proposition~\ref{prop:frame-free-rows}\,(iii)(b) is equivalent to the
conjugated row when $\ker W_t=0$, and otherwise implies it and in general
restricts further.

Declining the frame-free form costs nothing: the conjugated form is exact
on every support frame, tight or not (Theorem~\ref{thm:choi}), so imposing
it changes no value and forgoes only the redundant directions of
$\ker W_{\tau}$. That is why the condition may be resolved conservatively
rather than decided numerically.
\end{remark}
